% Ejemplo 1: ecuaciones de electrónica bien escritas
% Compilar: pdflatex 01_ecuaciones.tex  (dos veces para las referencias)
\documentclass[12pt,a4paper]{article}
\usepackage[utf8]{inputenc}
\usepackage[T1]{fontenc}
\usepackage[spanish]{babel}
\usepackage{amsmath,amssymb}
\usepackage{siunitx}
\sisetup{locale=FR, output-decimal-marker={,}}
\usepackage[hidelinks]{hyperref}

% Macros propias: escribe una vez, reutiliza siempre
\newcommand{\vv}[1]{\mathbf{#1}}
\DeclareMathOperator{\sen}{sen}

\title{Ecuaciones de electrónica con \LaTeX}
\author{Estudiante UNTELS}
\date{\today}

\begin{document}
\maketitle

\section{Modo en línea y modo destacado}
La ley de Ohm, $V = I R$, relaciona voltaje, corriente y resistencia.
Evita \texttt{\$\$...\$\$}; usa \texttt{\textbackslash[ ... \textbackslash]} o \texttt{equation}:
\begin{equation}
  P = VI = I^{2}R = \frac{V^{2}}{R}
  \label{eq:potencia}
\end{equation}
La ecuación~\eqref{eq:potencia} se cita con \verb|\eqref|, que añade los paréntesis.

\section{Ecuaciones alineadas}
Circuito RLC serie: la LVK y su forma en Laplace.
\begin{align}
  v(t) &= L\frac{di}{dt} + Ri + \frac{1}{C}\int_{0}^{t} i(\tau)\,d\tau \label{eq:rlc}\\
  V(s) &= \left(sL + R + \frac{1}{sC}\right) I(s) \label{eq:rlc_s}
\end{align}

\section{Matrices y sistemas}
Análisis de mallas:
\begin{equation}
  \begin{bmatrix}
    R_1 + R_2 & -R_2 \\
    -R_2 & R_2 + R_3
  \end{bmatrix}
  \begin{bmatrix} I_1 \\ I_2 \end{bmatrix}
  =
  \begin{bmatrix} V_1 \\ V_2 \end{bmatrix}
\end{equation}

Función definida por tramos (señal escalón con rampa):
\begin{equation}
  v(t) =
  \begin{cases}
    0, & t < 0 \\
    \dfrac{V_m}{T}\,t, & 0 \le t < T \\
    V_m, & t \ge T
  \end{cases}
\end{equation}

\section{Filtro pasa bajos de segundo orden}
\begin{equation}
  H(s) = \frac{\omega_0^{2}}{s^{2} + \dfrac{\omega_0}{Q}\,s + \omega_0^{2}},
  \qquad \omega_0 = \frac{1}{\sqrt{LC}}
\end{equation}

\section{Unidades con \texttt{siunitx}}
Con $R=\SI{4,7}{\kilo\ohm}$ y $C=\SI{100}{\nano\farad}$ la constante de tiempo es
$\tau = RC = \SI{470}{\micro\second}$ y la frecuencia de corte
$f_c = 1/(2\pi RC) \approx \SI{338,6}{\hertz}$.

\end{document}
